\documentclass{thesiscover}

\defaultfont

\begin{document}

\name{Ioan}{Dragomir}
\supervisor{Șl. Dr. Ing.}{Tassos}{Natsakis}

\thesistitle{Multi-channel biopotential DAQ with ROS2, IIO bindings}
\thesisproblem{O scurtă  descriere a temei proiectului de diplomă}
\thesiscontent{Presentation page, Declaration of authenticity, Summary, Contents, Introduction, Bibliographic research, Analysis design and implementation, Results and validation, Conclusion, Bibliography, Annexes}
\thesisplace{Technical University of Cluj-Napoca}
\thesisadvisors{}
\thesisemissiondate{02.07.2024}
\thesisdeliverydate{04.07.2024}

\idcard{RZ}{113316}
\CNP{5011113410095}
\examsession{2023-2024 September}
\declarationsigndate{02.07.2024}

\summaryrequirements{A 8/16 channel ECG and EMG DAQ with at least $1000\text{Hz}$ sample rate, and good signal characteristics ($\text{SNR} > 20\text{dB}$, high CMRR, PSRR). Both hardware, firmware and software components designed and implemented. Open source hardware and software.}
\summarysolutions{Hardware: MAX78000 microcontroller, AD4114 ADC, custom designed and built pre-amplifier and Driven Right Leg modules. Firmware implements the libIIO protocol, providing a serial iio context. Software provides lightweight bindings for libIIO, Python, and a ROS2 publisher node.}
\summaryresults{All signal quality objectives are checked, with room for improvement for the sample rate when using all 16 channels, and software capabilities.}
\summaryverification{The hardware components were tested individually using an ADALM2000-based test bench. The system as a whole was tested in various measurement scenarios: single/multiple channel EMG/ECG. The software components' throughput was tested to keep up with ADC output.}
\summarycontribution{Analog signal chain design. Firmware development and low level optimization of third party libraries. Thin software wrappers for interfacing with the rest of the system.}
\summarysources{M. Barbero, R. Merletti, and A. Rainoldi, "Atlas of Muscle Innervation Zones"; W. J. Esposito, "The ASPEN platform : Tools for mixed signal electronics, digital
signal processing, and biomedical electronics"}

\makecover
\end{document}
